\section{任意矩阵的奇异值分解}
\label{sec:03-Linear-Algebra/07-SVD-and-Data-Applications/01}

\subsection{特征向量为什么不够用}

我们花了很多篇幅研究特征向量：寻找那些被矩阵作用后只缩放、不转向的方向。\(Av=\lambda v\)——输入和输出落在同一条直线上。这个问题从一开始就把输入空间和输出空间当成了同一个空间，因此整个讨论只对方阵有意义。

但数据矩阵几乎从来不方。一个 \(100\times 30\) 的设计矩阵，把 30 维参数空间映射到 100 维观测空间。\(Av=\lambda v\) 要求等式两边向量的长度相等，而 \(Av\) 在 \(\R^{100}\) 中，\(v\) 在 \(\R^{30}\) 中——这个等式根本没有定义。

那么应该问一个什么问题来替代它？把问题的前提放宽：不再要求输入方向在映射后保持不变，只要求某些互相垂直的输入方向，在映射后的像也互相垂直。如果存在这样两组正交方向，矩阵在它们之间就只需做彼此独立的伸缩。

从几何上想象这个情景。把一个 \(2\times 2\) 矩阵作用在单位圆上，像通常是一个椭圆。椭圆有两条互相垂直的主轴，它们对应最长和最短的输出方向。而这两条主轴在原始单位圆上的原像，也是互相垂直的。高维也完全一样：单位球被压成一个可能在某方向退化（半轴长度为零）的椭球，椭球各半轴的方向和长度，正是 SVD 要找的左右奇异向量和奇异值。

"原像也互相垂直"目前只是几何观察。下面用 \(A^{\trans}A\) 把它变成严格推导。

\subsection{\texorpdfstring{从 \(A^{\trans} A\) 推出 SVD}{从 A^{\trans} A 推出 SVD}}

设 \(A\in\R^{m\times n}\)。直接问特征向量不行，但 \(A^{\trans}A\) 始终是 \(n\times n\) 的方阵：

\[
A^{\trans}A \in \R^{n\times n}.
\]

更重要的是，\(A^{\trans}A\) 是对称的——\((A^{\trans}A)^{\trans}=A^{\trans}A\)。而且它是半正定的：

\[
x^{\trans}A^{\trans}Ax = \|Ax\|_2^2 \ge 0.
\]

对称半正定矩阵的谱定理立刻生效：存在一组标准正交向量 \(v_1,\ldots,v_n\)，使得

\[
A^{\trans}Av_i = \lambda_i v_i,
\qquad
\lambda_1\ge\lambda_2\ge\cdots\ge\lambda_n\ge 0.
\]

\(\lambda_i\) 全部非负，所以可以定义奇异值

\[
\sigma_i = \sqrt{\lambda_i}.
\]

\(\sigma_i\) 就是 \(A\) 作用在 \(v_i\) 上时，输出向量的长度：\(\|Av_i\|^2=v_i^{\trans}A^{\trans}Av_i=\lambda_i=\sigma_i^2\)。

现在进入关键步骤。对于 \(\sigma_i>0\) 的那些方向（即 \(Av_i\ne 0\)），定义

\[
u_i = \frac{Av_i}{\sigma_i}.
\]

\(u_i\) 是单位向量。它们彼此正交吗？对 \(i\ne j\),

\[
u_i^{\trans}u_j
= \frac{v_i^{\trans}A^{\trans}Av_j}{\sigma_i\sigma_j}
= \frac{\lambda_j v_i^{\trans}v_j}{\sigma_i\sigma_j}
= \frac{\lambda_j\cdot 0}{\sigma_i\sigma_j}=0.
\]

推导只用了 \(v_i\) 的正交性——\(A^{\trans}A\) 是对称矩阵的谱定理直接提供了这个性质。因此我们得到了想要的结构：正交输入方向 \(\{v_i\}\) 被送到正交输出方向 \(\{u_i\}\)，缩放因子是 \(\sigma_i\)：

\[
Av_i = \sigma_i u_i.
\]

设 \(r=\rank(A)\)，即非零奇异值的个数。对于 \(i>r\)，\(\sigma_i=0\)，即 \(\|Av_i\|^2=0\)，所以 \(Av_i=0\)——这些方向被 \(A\) 完全压平。

现在构造矩阵 \(\sum_{i=1}^{r}\sigma_i u_i v_i^{\trans}\)。把它作用在任意 \(v_j\) 上：因为 \(v_i^{\trans}v_j=\delta_{ij}\)，只有 \(i=j\) 的项存活，结果是 \(\sigma_j u_j=Av_j\)。由于 \(\{v_1,\ldots,v_n\}\) 是 \(\R^n\) 的一组基，两个矩阵在一组基上的作用完全相同，因此两矩阵相等。这就得到紧致 SVD：

\[
A = U_r\Sigma_r V_r^{\trans}
= \sum_{i=1}^{r}\sigma_i u_i v_i^{\trans},
\]

其中 \(U_r\in\R^{m\times r}\) 的列是 \(u_1,\ldots,u_r\)（正交），\(V_r\in\R^{n\times r}\) 的列是 \(v_1,\ldots,v_r\)（正交），\(\Sigma_r=\diag(\sigma_1,\ldots,\sigma_r)\)。如果把两组向量分别补成 \(\R^m\) 和 \(\R^n\) 的完整正交基，就得到完整 SVD：\(A=U\Sigma V^{\trans}\)。

按从右向左的顺序读这个分解：\(V^{\trans}\) 把输入改写为"特殊输入方向"的坐标，\(\Sigma\) 对每个坐标独立伸缩（某些方向可能被压到零），\(U\) 把结果放到输出空间。正交变换 \(V^{\trans}\) 和 \(U\) 不改变向量长度，矩阵所有的放大、压缩和信息丢失，全部集中在 \(\Sigma\) 的对角线上。

\subsection{一个长方形矩阵的完整图景}

取最简单的长方形矩阵：

\[
A=
\begin{bmatrix}
2&0\\
0&1\\
0&0
\end{bmatrix}.
\]

它把 \(\R^2\) 嵌入 \(\R^3\)：第一坐标伸长两倍，第二坐标原样保留，第三坐标永远为 0。

计算 \(A^{\trans}A=\begin{bmatrix}4&0\\0&1\end{bmatrix}\)。特征向量显然是 \(v_1=e_1\)、\(v_2=e_2\)（标准基），对应奇异值 \(\sigma_1=2\)、\(\sigma_2=1\)。输出方向：\(u_1=Ae_1/2=(1,0,0)^{\trans}=e_1\)，\(u_2=Ae_2/1=(0,1,0)^{\trans}=e_2\)。第三个输出方向 \(e_3\) 永远不会被 \(A\) 产生——它属于 \(A^{\trans}\) 的零空间。

这个例子虽小，已经把四个基本子空间在 SVD 中的位置展示清楚：

\begin{itemize}
\tightlist
\item
  \(v_1,\ldots,v_r\) 张成 \(A\) 的行空间（输入中"有用"的方向）；其余右奇异向量张成 \(N(A)\)（输入中被压平的方向）；
\item
  \(u_1,\ldots,u_r\) 张成 \(A\) 的列空间（输出中可到达的方向）；其余左奇异向量张成 \(N(A^{\trans})\)（输出中永远到不了的方向）。
\end{itemize}

这些断言可以逐条验证。\(i>r\) 时 \(\|Av_i\|^2 = v_i^{\trans}A^{\trans}Av_i = 0\)，故靠后的右奇异向量确实在 \(N(A)\) 中。而 \(A=\sum_{i\le r}\sigma_i u_i v_i^{\trans}\) 表明：每个输出都是前 \(r\) 个 \(u_i\) 的组合，故列空间恰由它们张成；每一行都是前 \(r\) 个 \(v_i\) 的组合，故行空间恰由它们张成。

因此 SVD 不只是一个计算公式。它同时给出了输入中被保留的方向、被消灭的方向，以及输出中可以到达和永远到不了的方向。一个线性映射完整的信息结构，全部陈列在两组正交向量和一条对角线上。

再看一个秩一矩阵：

\[
B=\begin{bmatrix}1&1\\1&1\end{bmatrix}.
\]

沿 \(v_1=(1,1)^{\trans}/\sqrt{2}\) 的输入被放大两倍并送到同一方向；沿 \(v_2=(1,-1)^{\trans}/\sqrt{2}\) 的输入被完全压成零。奇异值：\(\sigma_1=2\)，\(\sigma_2=0\)。

零奇异值不是一个不幸的计算误差。它告诉你一件结构上的事：输入的"差异方向"一旦被抹掉，就无法从输出中恢复。如果你试图求 \(B^{-1}\)，算法会崩溃——不是因为实现有 bug，而是因为你要求的事情在数学上就不可能。

\subsection{需要保留的边界}

奇异值是唯一的，但奇异向量并不总是唯一。每一对 \((u_i,v_i)\) 可以同时乘以 \(-1\) 而不改变 SVD。如果某个奇异值重复（比如 \(\sigma_2=\sigma_3\)），对应子空间内的任意正交基都同样合法。因此，当你试图解释数据中的"第几个奇异向量"时，先确认对应的奇异值是否与邻近值充分分离。两个几乎相等的奇异值意味着对应方向尚未被数据可靠地定下来——任何正交旋转都是近似等价的。

还要区分奇异值与特征值。奇异值总是非负实数，适用于任意矩形矩阵。特征值可以为负、为复数，描述方阵的不变方向，且只在可对角化（或至少能走 Jordan 形）时才有完整的谱刻画。只有对实对称（或更一般地对正规矩阵的奇异值），才有简单关系：奇异值等于特征值的绝对值。

SVD 对所有矩阵无条件存在。这是它与特征分解的根本区别。下一节会把 SVD 用于求解看似不可解的线性方程组——当逆矩阵不存在时，SVD 给出了一个信息损失最小的替代方案。

\subsection{延伸与来源}

\begin{itemize}
\tightlist
\item
  MIT 18.06SC：Singular Value Decomposition
\item
  MIT 18.065：Matrix Methods in Data Analysis, Signal Processing, and Machine Learning
\end{itemize}
